\chapter{Discapacidad Auditiva y Comunidad Sorda}
\label{cp:discapacidad-auditiva}

\insertminitoc
\parindent0pt

Para comprender el contexto en que se desarrolla esta investigación, es necesario examinar
la discapacidad auditiva desde dos perspectivas complementarias: la médica, que permite
entender sus causas, clasificación y prevalencia, y la sociocultural, que reconoce a las
personas sordas como miembros de una comunidad lingüística con identidad propia. Este
capítulo aborda ambas perspectivas, con especial atención a la situación de los niños
sordos, las barreras de comunicación que enfrentan con la población oyente y el contexto
específico de Honduras donde se enmarca la investigación.

\section{Discapacidad Auditiva: Definición, Clasificación y Prevalencia}

La pérdida auditiva se define como la disminución parcial o total de la capacidad de
percibir sonidos. La Organización Mundial de la Salud (\acrshort{oms}) clasifica la pérdida auditiva
en grados según el umbral de audición medido en decibelios (dB), evaluado mediante el
promedio de tonos puros en las frecuencias de 500, 1000, 2000 y 4000 Hz. La
Tabla~\ref{tab:clasificacion-perdida-auditiva} resume la clasificación vigente, que
distingue siete grados que van desde la audición normal hasta la pérdida completa.

\begin{table}[h!]
\centering
\begin{tabular}{lc}
\hline
\textbf{Grado de pérdida auditiva} & \textbf{Umbral de audición (dB)} \\
\hline
Audición normal & -10.0 a 19.9 \\
Pérdida leve & 20.0 a 34.9 \\
Pérdida moderada & 35.0 a 49.9 \\
Pérdida moderadamente severa & 50.0 a 64.9 \\
Pérdida severa & 65.0 a 79.9 \\
Pérdida profunda & 80.0 a 94.9 \\
Pérdida completa & 95.0 o superior \\
\hline
\end{tabular}
\caption[Clasificación de la pérdida auditiva según la OMS]{Clasificación de la pérdida auditiva según la \acrshort{oms}, por grado y umbral de audición. Fuente: elaboración propia a partir de \cite{olusanya_hearing_2019}.}
\label{tab:clasificacion-perdida-auditiva}
\end{table}

Como se observa en la Tabla~\ref{tab:clasificacion-perdida-auditiva}, cada grado
corresponde a un rango del umbral auditivo, de manera que a mayor umbral en decibelios
mayor es la pérdida. Sobre la base de esta escala, la \acrshort{oms} considera como pérdida
auditiva discapacitante aquella superior a 40 dB en adultos y superior a 30 dB en niños \cite{who_deafness_2024}.

En términos de prevalencia, el Informe Mundial sobre la Audición publicado por la \acrshort{oms}
en 2021 estima que más de 1\,500 millones de personas en el mundo presentan algún grado
de pérdida auditiva, de las cuales aproximadamente 430 millones requieren atención
especializada \cite{who_world_report_2021}. De este total, se estima que 34 millones son niños \cite{who_deafness_2024}. El informe proyecta
que, sin intervenciones preventivas adecuadas, para el año 2050 cerca de 2\,500 millones de
personas vivirán con algún grado de pérdida auditiva y al menos 700 millones necesitarán
servicios de rehabilitación \cite{who_world_report_2021}.

La pérdida auditiva puede ser de origen congénito o adquirido. Entre las causas congénitas
se encuentran factores genéticos, complicaciones durante el embarazo y el parto, bajo peso
al nacer y la ictericia neonatal. Entre las causas adquiridas destacan las infecciones crónicas
del oído, la exposición prolongada a ruidos intensos, el uso de medicamentos ototóxicos y
el envejecimiento \cite{who_deafness_2024}. En el caso de los niños, aproximadamente 200 millones padecen
infecciones crónicas del oído que pueden prevenirse o tratarse, lo que subraya la
importancia de la intervención temprana \cite{who_world_report_2021}.

\section{La Comunidad Sorda y la Cultura Visual}

Las personas sordas conforman una comunidad con una identidad cultural y lingüística
propia que trasciende la condición audiológica. La Federación Mundial de Sordos (\acrshort{wfd}),
organización internacional que representa aproximadamente a 70 millones de personas
sordas en 135 países, ha sostenido de forma consistente que las personas sordas no se
definen por una condición médica deficitaria, sino por su pertenencia a una minoría
lingüística y cultural cuyo rasgo central es el uso de una lengua de señas nacional \cite{wfd_cultural_2018}.

La Convención sobre los Derechos de las Personas con Discapacidad (\acrshort{cdpd}) de las
Naciones Unidas reconoce en su Artículo 30 que las personas con discapacidad tienen
derecho al reconocimiento y apoyo de su identidad cultural y lingüística específica,
incluyendo las lenguas de señas y la cultura sorda \cite{onu_cdpd_2006}. La \acrshort{wfd} ha señalado que este
reconocimiento implica una interseccionalidad singular: las personas sordas pertenecen
simultáneamente al ámbito de la discapacidad y al de las minorías lingüísticas, una
condición que no comparten otros grupos \cite{wfd_sign_language_rights_2024}.

La comunicación en la comunidad sorda se organiza en torno a la modalidad visual y
gestual. A diferencia de las lenguas orales, que operan en el canal auditivo-vocal, las
lenguas de señas utilizan el canal visual-gestual, lo que implica una organización espacial
del lenguaje. Esta diferencia no representa una limitación, sino una modalidad alternativa
igualmente compleja y expresiva. Las lenguas de señas nacionales se originan en las
comunidades sordas y se desarrollan con ellas, por lo que cada país o región puede tener
su propia lengua de señas, con gramática y vocabulario independientes de la lengua oral
del mismo territorio \cite{wfd_primacy_2024}.

\section{La Infancia Sorda y el Acceso Temprano a la Lengua de Señas}

La adquisición del lenguaje en los primeros años de vida constituye un proceso
determinante para el desarrollo cognitivo, emocional y social de cualquier niño. Humphries
et al. documentaron que los niños adquieren el lenguaje sin instrucción formal siempre que
estén expuestos de manera regular y significativa a una lengua humana accesible. Sin
embargo, debido a los cambios en la plasticidad cerebral durante la primera infancia, los
niños que no adquieren una primera lengua en los primeros años podrían no alcanzar nunca
una competencia lingüística plena en ningún idioma \cite{humphries_language_2012}.

Este fenómeno, conocido como deprivación lingüística, tiene consecuencias directas sobre
el desarrollo de actividades cognitivas que dependen de una base lingüística sólida, como
la lectoescritura, la organización de la memoria y la manipulación numérica \cite{humphries_language_2012}. La
adquisición de una lengua de señas está sujeta a las mismas restricciones temporales que
la adquisición de una lengua oral: si el niño sordo no accede a una lengua de señas durante
el período crítico, las consecuencias sobre su desarrollo lingüístico pueden ser permanentes.

La \acrshort{wfd} ha enfatizado en su documento de posición sobre los derechos lingüísticos de los
niños sordos que la necesidad de adquisición natural del lenguaje a través de una lengua
de señas es un derecho fundamental de todos los niños sordos, y que los métodos
educativos que mejor promueven el desarrollo de su identidad cultural y lingüística son los
modelos bilingües con lengua de señas completa \cite{wfd_language_rights_2016}.

A pesar de la evidencia científica, Humphries et al. señalaron que el 80\,\% de los niños
sordos nacen en familias oyentes, lo que con frecuencia retrasa o impide el acceso temprano
a una lengua de señas \cite{humphries_language_2012}. Esta situación es especialmente grave en países en desarrollo
donde los servicios de detección temprana, los programas de intervención y la disponibilidad
de profesionales en lengua de señas son limitados o inexistentes.

\section{Barreras de Comunicación con Personas Oyentes}

La barrera fundamental entre la comunidad sorda y la oyente radica en la asimetría
lingüística: mientras que las personas sordas necesitan la lengua de señas como medio
principal de comunicación, la mayoría de la población oyente no la conoce. Esta asimetría
genera exclusión en múltiples ámbitos de la vida cotidiana.

En el entorno educativo, la escasez de intérpretes de lengua de señas y de docentes con
competencia en la lengua constituye una de las barreras más documentadas. La \acrshort{cdpd}
establece en su Artículo 24 la obligación de los Estados de asegurar que la educación de
las personas sordas se imparta en las lenguas y los modos y medios de comunicación más
apropiados, y en entornos que permitan alcanzar su máximo desarrollo académico y
social \cite{onu_cdpd_2006}.

En el entorno familiar, la barrera se origina desde el nacimiento cuando el niño sordo nace
en una familia oyente que no conoce la lengua de señas. La comunicación temprana se ve
restringida a gestos naturales y estrategias ad hoc que no constituyen un sistema
lingüístico completo, lo que afecta el vínculo afectivo y la transmisión cultural entre padres
e hijos \cite{humphries_language_2012}.

En el entorno comunitario y de servicios, las personas sordas enfrentan dificultades para
acceder a servicios básicos de salud, trámites legales y atención de emergencias. La
ausencia de mecanismos de comunicación accesibles obliga a depender de familiares como
intermediarios, lo que compromete la privacidad y la autonomía de la persona sorda. La
tecnología móvil ha sido señalada como una vía prometedora para reducir estas barreras,
siempre que las soluciones estén adaptadas a la lengua de señas local y sean accesibles
económicamente \cite{who_world_report_2021}.

\section{Situación de las Personas Sordas en Honduras}

Luego de revisar el panorama general de la discapacidad auditiva y sus implicaciones
comunicativas y educativas, es necesario situar esa realidad en el contexto hondureño,
que es donde se desarrolla y hacia donde se dirigen los resultados de esta investigación.

En Honduras, la Comisión Nacional de los Derechos Humanos (\acrshort{conadeh}) ha reportado
que existen aproximadamente 100\,000 personas con discapacidad auditiva en el país. Sin
embargo, únicamente alrededor del 1\,\% de esta población cuenta con un carné oficial de
identificación como persona con discapacidad, lo que evidencia un bajo nivel de
reconocimiento institucional y un acceso limitado a los servicios y beneficios a los que
tienen derecho \cite{conadeh_discriminados}.

Williams y Parks, en una encuesta sociolingüística realizada para SIL International,
documentaron que la comunidad sorda hondureña se comunica a través del Lenguaje de
Señas Hondureño (\acrshort{lesho}) y que más de 3\,000 personas han recibido formación en esta
lengua a través de cursos ofrecidos por la Asociación de Sordos de Honduras (\acrshort{ash}) y la
Universidad Nacional Autónoma de Honduras (\acrshort{unah}). El estudio también registró
variaciones regionales en el \acrshort{lesho} y señaló que la mayoría de las personas sordas
aprenden la lengua a través de su escuela o de las clases impartidas por la \acrshort{ash} \cite{williams_encuesta_2015}.

En el ámbito educativo, la \acrshort{unah} ha ofrecido cursos de \acrshort{lesho} de forma ininterrumpida
durante más de 23 años, y en 2024 el Programa de Servicios a Estudiantes con Necesidades
Especiales (\acrshort{prosene}) lanzó el Diccionario Digital de \acrshort{lesho} con más de 700 señas en
formato de video, orientado inicialmente al vocabulario del entorno universitario. El
diccionario fue desarrollado por personal y estudiantes sordos adscritos al \acrshort{prosene} y
constituye el primer recurso digital de su tipo para esta lengua \cite{unah_diccionario_2024}.

A pesar de estos avances, el acceso a intérpretes certificados fuera del ámbito educativo
universitario sigue siendo escaso, especialmente en zonas rurales y ciudades intermedias.
No existen a la fecha herramientas tecnológicas locales que permitan la comunicación
autónoma en tiempo real entre personas sordas y oyentes sin la intermediación de un
intérprete humano.
